% =============================================================================
% БИЛЕТ 17
% =============================================================================

\ticket{17}{}

\question{Дифференцируемость композиции дифференцируемых отображений}

Для начала заботаем, что такое отображение, это функция между пространствами, которая из $n$-мерного пространства в $m$-мерное. Заботали, далее надо понять что матрица этих отображений(матрица Якоби) при композиции отображений перемножается. И наконец, если оба отображения дифференцируемые, то их композиция тоже дифференцируемая. Отсюда у нас есть формула для производной композиции, которая говорит, что производная композиции равна произведению производных.
\begin{theorem}[Дифференцируемость композиции]
Пусть $f: U \subseteq \mathbb{R}^n \to \mathbb{R}^m$ и $g: V \subseteq \mathbb{R}^m \to \mathbb{R}^p$ - дифференцируемые отображения, где $U$ и $V$ - открытые множества, и $f(U) \subseteq V$. Тогда композиция $h = g \circ f: U \to \mathbb{R}^p$ также дифференцируема, и её производная в точке $x \in U$ задаётся формулой
\[Dh(x) = Dg(f(x)) \cdot Df(x),\]
где $Dh(x)$, $Dg(f(x))$ и $Df(x)$ - матрицы Якоби отображений $h$, $g$ и $f$ соответственно.  
\end{theorem}

\begin{proof}
Пусть $f: U \subseteq \mathbb{R}^n \to \mathbb{R}^m$ и $g: V \subseteq \mathbb{R}^m \to \mathbb{R}^p$ - дифференцируемые отображения, где $U$ и $V$ - открытые множества, и $f(U) \subseteq V$. Нам нужно показать, что композиция $h = g \circ f: U \to \mathbb{R}^p$ также дифференцируема и найти её производную.
Пусть $x \in U$ и рассмотрим $h(x) = g(f(x))$. Поскольку $f$ дифференцируемо в точке $x$, существует линейное отображение $Df(x): \mathbb{R}^n \to \mathbb{R}^m$, такое что
\[f(x + h) = f(x) + Df(x)h + o(\|h\|)\]
для всех $h$ в окрестности нуля. Аналогично, поскольку $g$ дифференцируемо в точке $f(x)$, существует линейное отображение $Dg(f(x)): \mathbb{R}^m \to \mathbb{R}^p$, такое что
\[g(y + k) = g(y) + Dg(f(x))k + o(\|k\|)\]
для всех $k$ в окрестности нуля, где $y = f(x)$.
Теперь рассмотрим $h(x + h) = g(f(x + h))$. Подставляя разложение для $f(x + h)$, получаем
\[h(x + h) = g(f(x) + Df(x)h + o(\|h\|)).\]
Используя разложение для $g$, получаем
\[h(x + h) = g(f(x)) + Dg(f(x))(Df(x)h + o(\|h\|)) + o(\|Df(x)h + o(\|h\|)\|).\]  и из этого следует, что
\[h(x + h) = g(f(x)) + Dg(f(x))Df(x)h + o(\|h\|).\]
Таким образом, мы видим, что $h$ дифференцируемо в точке $x$, и его производная задаётся формулой
\[Dh(x) = Dg(f(x)) \cdot Df(x).\]
\end{proof}

\question{Критерий Лебега интеграла Римана}


% Ответ...

